\section{Three critical points on the fixed submanifold}
\label{sec:three-points}

We next construct a function on the fixed submanifold with three
critical points. The comparison with functions on the ambient quadric
will follow in Section~\ref{sec:critical-bound}.
Three-critical-point functions on products of spheres belong to the
classical study of critical-point minima and category by
Takens~\cite{Takens1968}; the $S^2\times S^2$ case is recorded by
Mare~\cite[Proposition~3.2.5]{Mare1997}. We use the explicit formula below
and determine its entire critical set directly, without importing an
existence theorem or assuming that its critical points are Morse.

\begin{proposition}
\label{prop:three-points}
Let $p=(0,0,1)$ and $e=(1,0,0)$ in $\R^3$.
The smooth function
\begin{equation}
\label{eq:three-point-function}
  f\colon S^2\times S^2\longrightarrow\R,
  \qquad f(x,y)=e\cdot x+(x-p)\cdot y,
\end{equation}
has exactly three critical points.
\end{proposition}

\begin{proof}
For a fixed $x\ne p$, the $y$-dependence is a height function on the
second sphere, with critical points $y=\pm(x-p)/|x-p|$.
At $x=p$ that height function is constant, but criticality of the full
function still requires its derivative in the $x$ directions to vanish.
We therefore separate this exceptional fiber from the two height branches.

The constrained critical-point equations are
\begin{equation}
\label{eq:critical-multipliers}
  e+y=\lambda x,\qquad x-p=\mu y,
  \qquad |x|=|y|=1,
\end{equation}
for real multipliers $\lambda,\mu$.
If $x=p$, the first equation gives $y=-e+\lambda p$.
Its unit norm forces $\lambda=0$, yielding the unique point
\[
  q_0=(p,-e),\qquad f(q_0)=0.
\]

Suppose now that $x\ne p$ and set $t=|x-p|>0$.
The second equation in \eqref{eq:critical-multipliers} gives
$y=s(x-p)/t$ for a sign $s\in\{1,-1\}$.
The first then becomes
\[
  t e-s p=(\lambda t-s)x.
\]
Since the left side is nonzero, $\lambda t-s\ne0$ and $x$ lies in
the plane spanned by $e$ and $p$. Write $x=(u,0,w)$.
The definition of $t$ gives $w=1-t^2/2$.
Taking the first equation in \eqref{eq:critical-multipliers}
against the tangent vector $(w,0,-u)$ gives
\[
  w+\frac{s u}{t}=0,
  \qquad\text{so}\qquad u=-s t w.
\]
Consequently the unit-norm condition reduces to
\begin{equation}
\label{eq:critical-factorization}
  0=(1-t^2/2)^2(1+t^2)-1
   =\frac{t^4(t^2-3)}4.
\end{equation}
Because $t>0$, this forces $t=\sqrt3$, $w=-1/2$, and
$u=s\sqrt3/2$. In particular, the other pole $x=-p$, for which $t=2$,
is excluded. The two resulting candidates are
\begin{equation}
\label{eq:other-critical-points}
  q_s=
  \left(
    \left(\frac{s\sqrt3}{2},0,-\frac12\right),
    \left(\frac12,0,-\frac{s\sqrt3}{2}\right)
  \right),
  \qquad s\in\{1,-1\}.
\end{equation}
Both satisfy \eqref{eq:critical-multipliers} with
$\lambda=\mu=s\sqrt3$, and
\[
  f(q_s)=s\,\frac{3\sqrt3}{2}.
\]
The cases $x=p$ and $x\ne p$ exhaust the domain, proving the count.
\end{proof}

\begin{remark}
\label{rem:degenerate-critical-point}
The critical point $q_0$ is degenerate. Indeed, in local coordinates
\[
  x=(u,v,\sqrt{1-u^2-v^2}),\qquad
  y=(-\sqrt{1-a^2-b^2},a,b),
\]
the function has expansion
\[
  f=va+O\bigl(\lVert(u,v,a,b)\rVert^3\bigr).
\]
Its Hessian therefore has eigenvalues $1,-1,0,0$, so its rank is two.
The complete critical-point calculation proves that this degenerate point
is isolated. This is why a count restricted to Morse functions would not
capture the example.
\end{remark}
